% GENERATED FILE. Edit canonical lessons, then run npm run generate:books.
% canonical-source-hash: fnv1a64:8b602d220c8fc709
% canonical-lessons: ES-C17-condicional, ES-C17-comer-condicional, ES-C17-vivir-condicional

\chapter{The Conditional --- Saying What Would Happen}
\label{ch:spanish-condicional}
\hlchaptermodality{80}

\begin{chapteropening}
\textbf{By the end of this chapter:} \emph{I can say what I would do, in the singular, across all three verb families.}

The same infinitive, a different auxiliary tense on the end.
\end{chapteropening}

\section[hablaría]{hablaría — the whole infinitive meets -ía}

\label{lesson:ES-C17-condicional}



\begin{quote}
\textbf{Hablaré español} places speaking in the future. \textbf{Hablaría español} means “I would speak Spanish” in an imagined or conditional situation. Today the situation can remain unstated; learn the form and its core “would” meaning first.
\end{quote}

\begin{grammarlens}[title={Grammar Lens: three familiar persons}]
Keep \textbf{hablar} whole, then add the learned singular conditional endings:

\noindent
\begin{tabularx}{\linewidth}{@{}>{\raggedright\arraybackslash}X>{\raggedright\arraybackslash}X@{}}
\toprule
\textbf{person} & \textbf{conditional} \\
\midrule
yo & hablar\textbf{ía} \\
tú & hablar\textbf{ías} \\
él / ella / usted & hablar\textbf{ía} \\
\bottomrule
\end{tabularx}

Say \textbf{hablaría · hablarías · hablaría}. The first and third forms match. The accent keeps \textbf{í-a} in separate syllables, just as it did in \textbf{comía}. Plural persons wait.
\end{grammarlens}

\begin{cousinweb}
The conditional developed from an infinitive beside past or imperfect forms of the same “have” verb that helped form the Romance future. That is why the whole infinitive remains visible again and why the endings resemble the \textbf{-ía} row. The history passes through older Romance forms; modern \textbf{hablaría} is not made by attaching the complete modern word \textbf{había}.
\end{cousinweb}

\subsection*{Your turn}
Say these aloud:

\begin{itemize}
\raggedright
  \item “hablaría, hablarías, hablaría”
  \item “Hablaré español. Hablaría español.”
  \item “ha-bla-rí-a”
\end{itemize}

\begin{tcolorbox}[breakable,colback=teal!4,colframe=teal!35!black,title={Before you move on}]
Give the three learned forms. (\emph{Hablaría, hablarías, hablaría.}) What is the core meaning? (“Would speak.”) Which older construction stands behind the row? (An infinitive plus past forms of \textbf{haber}.) Next: give \textbf{comer} the same ending set.
\end{tcolorbox}
\section[comería]{comería — keep -er before “would”}

\label{lesson:ES-C17-comer-condicional}



\begin{quote}
\textbf{Comeré} means “I will eat.” Keep \textbf{comer} whole and exchange the future ending for \textbf{-ía}: \textbf{comería} — “I would eat.”
\end{quote}

\begin{grammarlens}[title={Grammar Lens: the singular -er row}]
\noindent
\begin{tabularx}{\linewidth}{@{}>{\raggedright\arraybackslash}X>{\raggedright\arraybackslash}X@{}}
\toprule
\textbf{person} & \textbf{conditional} \\
\midrule
yo & comer\textbf{ía} \\
tú & comer\textbf{ías} \\
él / ella / usted & comer\textbf{ía} \\
\bottomrule
\end{tabularx}

Say \textbf{comería · comerías · comería}. Compare \textbf{hablaría / comería}: one bounded ending set serves both complete infinitives. Plural persons wait.
\end{grammarlens}

\subsection*{Your turn}
Say these aloud:

\begin{itemize}
\raggedright
  \item “comería, comerías, comería”
  \item “comeré, comería”
  \item “hablaría, comería”
\end{itemize}

\begin{tcolorbox}[breakable,colback=teal!4,colframe=teal!35!black,title={Before you move on}]
Give the three learned forms. (\emph{Comería, comerías, comería.}) Which piece stays whole? (\textbf{Comer}.) Which endings match \textbf{hablaría}? (\textbf{-ía, -ías, -ía}.) Next: complete the regular pattern with \textbf{vivir}.
\end{tcolorbox}
\section[viviría]{viviría — the regular conditional closes its small frame}

\label{lesson:ES-C17-vivir-condicional}



\begin{quote}
You can move from \textbf{comeré} to \textbf{comería}. Apply the same move to \textbf{viviré}: \textbf{viviría} — “I would live.”
\end{quote}

\begin{grammarlens}[title={Grammar Lens: the singular -ir row}]
\noindent
\begin{tabularx}{\linewidth}{@{}>{\raggedright\arraybackslash}X>{\raggedright\arraybackslash}X@{}}
\toprule
\textbf{person} & \textbf{conditional} \\
\midrule
yo & vivir\textbf{ía} \\
tú & vivir\textbf{ías} \\
él / ella / usted & vivir\textbf{ía} \\
\bottomrule
\end{tabularx}

Say \textbf{viviría · vivirías · viviría}. The complete regular frame is now \textbf{hablaría, comería, viviría}. Plural persons wait.
\end{grammarlens}

\begin{grammarlens}[title={Grammar Lens: meaning with known words}]
\textbf{Viviría en Madrid} says “I would live in Madrid.” The lesson adds no new place or preposition.
\end{grammarlens}

\subsection*{Your turn}
Say these aloud:

\begin{itemize}
\raggedright
  \item “viviría, vivirías, viviría”
  \item “hablaría, comería, viviría”
  \item “Viviría en Madrid.”
\end{itemize}

\begin{tcolorbox}[breakable,colback=teal!4,colframe=teal!35!black,title={Before you move on}]
Give the three \textbf{vivir} forms. (\emph{Viviría, vivirías, viviría.}) Which regular verb families share the ending row? (All three.) Next: three already-known verbs change the stem but keep both sets of endings.
\end{tcolorbox}
